\section*{الفصل \ref{ch:g12:metamorphic-rocks} --- الصخور المتحولة}

\begin{solution}{exo:g12:metamorphic-rocks:1}
(أ) \omterm{def:g12:metamorphic-rocks:polymorph}{متعددات الأشكال}؛ (ب) \omterm{def:g12:metamorphic-rocks:index-mineral}{خط تساوي الدرجة}؛ (ج) \omterm{def:g12:metamorphic-rocks:contact-regional}{هالة التحول}؛ (د) \omterm{def:g12:metamorphic-rocks:foliation}{التورق}؛
(ه) \omterm{def:g12:metamorphic-rocks:lithostatic}{الضغط الصخري}.
\end{solution}

\begin{solution}{exo:g12:metamorphic-rocks:2}
(1) الرخام؛ (2) الكوارتزيت؛ (3) الأردواز؛ (4) النيس؛ (5) الهورنفلس.
\end{solution}

\begin{solution}{exo:g12:metamorphic-rocks:3}
$P = \rho g z$: على عمق \qty{10}{km}، $2830 \times 9.81 \times \num{10000} \approx
\qty{2.78e8}{Pa} = \qty{0.278}{GPa}$؛ وعلى عمق \qty{20}{km}، ضعف ذلك،
\qty{0.555}{GPa}.
\end{solution}

\begin{solution}{exo:g12:metamorphic-rocks:4}
(1) الكيانيت (فخط الكيانيت--الأندالوسيت قريب من \qty{0.25}{GPa} عند
\qty{400}{\celsius}). (2) الأندالوسيت. (3) السيليمانيت.
\end{solution}

\begin{solution}{exo:g12:metamorphic-rocks:5}
(1) الأردواز، فالفيليت، فالشيست، فالنيس. (2) في الشيست. (3) النيس.
\end{solution}

\begin{solution}{exo:g12:metamorphic-rocks:6}
(1) $2830 \times 9.81 \times \num{35000} \approx \qty{9.72e8}{Pa}$، أي
$\num{9.72e8}/101325 \approx \num{9600}$ ضغط جوي. (2) $1025 \times 9.81
\times 4000 \approx \qty{4.02e7}{Pa} = \qty{40.2}{MPa}$. (3) $972/40.2
\approx 24$ مرة.
\end{solution}

\begin{solution}{exo:g12:metamorphic-rocks:7}
(1) $z = \num{0.375e9}/(2830 \times 9.81) \approx \qty{13.5}{km}$.
(2) $(504 - 10)/13.5 \approx \qty{36.6}{\celsius}$ لكل كيلومتر. (3) لا:
فالصخر يجب أن يقع تمامًا عند نقطة واحدة من المخطط؛ ومثل هذه الصخور
نادرة.
\end{solution}

\begin{solution}{exo:g12:metamorphic-rocks:8}
(1) $23/3 \approx \qty{8}{mm}$. (2) نحو 8 ملايين سنة. (3) بقي الصخر
في العمق، عند درجة حرارة عالية، ملايين السنين.
\end{solution}

\begin{solution}{exo:g12:metamorphic-rocks:9}
(1) أربعة: \omterm{def:g12:metamorphic-rocks:index-mineral}{خطوط تساوي درجة} العقيق والستوروليت والكيانيت والسيليمانيت، بهذا
الترتيب. (2) نحو الغرب. (3) عند الطرف الغربي، في نطاق السيليمانيت. فقد
عبرت صخوره خط الكيانيت--السيليمانيت: وإذ دُفنت على عمق نحو \qty{20}{km}
(\qty{0.56}{GPa})، كانت أسخن من نحو \qty{590}{\celsius}.
\end{solution}

\begin{solution}{exo:g12:metamorphic-rocks:10}
(1) الوحل يصير هورنفلسًا، والحجر الجيري رخامًا، والحجر الرملي
كوارتزيتًا. (2) لقد سُخّنت دون أن تُعصر: فلم يرتب ضغط موجَّه
معادنها. (3) الصخور المحيطة به ساخنة أصلًا، فلا تزيدها
الصهارة إلا قليلًا من الحرارة.
\end{solution}

\begin{solution}{exo:g12:metamorphic-rocks:11}
(1) متعامدة على الدفع: أوجهها المسطحة تواجه الشرق والغرب، فتقف الرقائق
في مستويات شمالية--جنوبية. (2) رأسي، يمتد من الشمال إلى الجنوب.
(3) حبيبات الكالسيت والكوارتز ليست مسطحة ولا طويلة، فحتى حين
تتبلور من جديد تحت العصر لا يوجد ما يصطف.
\end{solution}

\begin{solution}{exo:g12:metamorphic-rocks:12}
(1) $z = (T - 10)/30$، إذن $P = 0.0278\,(T - 10)/30 \approx
0.000927\,(T - 10)$. (2) $0.000927\,(T - 10) = 0.577 + 0.00209\,(T - 600)$
تعطي $0.001163\,T \approx 0.668$، إذن $T \approx \qty{574}{\celsius}$، و$P
\approx 0.000927 \times 564 \approx \qty{0.523}{GPa}$ و$z = (574 - 10)/30
\approx \qty{18.8}{km}$. (3) الكيانيت فوق ذلك العمق، والسيليمانيت تحته.
\end{solution}

\begin{solution}{exo:g12:metamorphic-rocks:13}
(1) $10 + 10 \times 35 = \qty{360}{\celsius}$. (2) $10 + 30 \times 35 =
\qty{1060}{\celsius}$. (3) الضغط على عمق \qty{35}{km} نحو
\qty{0.97}{GPa}؛ وعند \qty{360}{\celsius} تقع تلك النقطة في عمق حقل
الكيانيت. أما النقطة الثانية فتقع بعيدًا إلى يمين المخطط، في حقل
السيليمانيت، عند درجات حرارة تبدأ فيها صخور القشرة بالانصهار.
(4) يحتاجان إلى ضغط عالٍ عند درجة حرارة منخفضة، وهذا اجتماع لا
توفره إلا القشرة الباردة المحمولة إلى الأسفل في \omterm{def:g11:subduction-zones:subduction}{نطاق اندساس}.
\end{solution}

\begin{solution}{exo:g12:metamorphic-rocks:14}
(1) الكيانيت $1000/3.61 \approx \qty{277}{cm^3}$؛ والسيليمانيت $1000/3.24
\approx \qty{309}{cm^3}$؛ والأندالوسيت $1000/3.15 \approx \qty{317}{cm^3}$.
(2) $(317 - 277)/317 \approx 0.127$، أي نحو \qty{13}{\percent}. (3) تحت
الضغط العالي يُحابى الترتيب الأكثر تراصًا للذرات؛ والكيانيت
يرصّ الذرات نفسها في أقل حيز.
\end{solution}

\begin{solution}{exo:g12:metamorphic-rocks:15}
(1) \qty{0.555}{GPa} على عمق \qty{20}{km}، و\qty{0.694}{GPa} على عمق \qty{25}{km}.
(2) $T = 600 + (P - 0.577)/0.00209$: نحو \qty{590}{\celsius} عند
\qty{0.555}{GPa} و\qty{656}{\celsius} عند \qty{0.694}{GPa}. وعند
\qty{700}{\celsius} كانت الصخور أسخن من كلتيهما، في حقل
السيليمانيت: فهذا متفق. (3) $(700 - 10)/20 \approx \qty{34.5}{\celsius}$
و$(700 - 10)/25 \approx \qty{27.6}{\celsius}$ لكل كيلومتر، وهما قريبان من
التدرج النموذجي. (4) من $\qty{20}{km}$ إلى $\qty{25}{km}$ في 40 مليون سنة:
من 0.50 إلى \qty{0.63}{km} في كل مليون سنة، أي من 0.50 إلى \qty{0.63}{mm} في
السنة.
\end{solution}

\begin{solution}{pb:g12:metamorphic-rocks:1}
\textbf{1.} إنه متورق، والميكا فيه كبيرة بما يكفي لتُرى.
والميكا، المسكوفيت والبيوتيت، هي التي تصنع \omterm{def:g12:metamorphic-rocks:foliation}{الشستوزية}.

\textbf{2.} \ce{Al2SiO5}؛ إنها \omterm{def:g12:metamorphic-rocks:polymorph}{متعددات أشكال}.

\textbf{3.} حجما الكيلوغرام في النسبة $3.15/3.61 \approx
0.873$: فالكيانيت أصغر بنحو \qty{13}{\percent}. والضغط يحابي
الكيانيت.

\textbf{4.} نطاق الكيانيت. وقد عبر في طريق نزوله \omterm{def:g12:metamorphic-rocks:index-mineral}{خطوط تساوي درجة} البيوتيت،
والعقيق، والستوروليت، والكيانيت.

\textbf{5.} نحو 8 ملايين سنة.

\textbf{6.} $P = 0.577 + 0.00209 \times 0 = \qty{0.577}{GPa}$.

\textbf{7.} \qty{0.70}{GPa} فوق \qty{0.577}{GPa}، في جانب الكيانيت من
خط الكيانيت--السيليمانيت؛ وهي أيضًا أعلى بكثير من حقل الأندالوسيت،
الذي يقع عند \qty{600}{\celsius} تحت نحو \qty{0.23}{GPa}.

\textbf{8.} \qty{0.577}{GPa}.

\textbf{9.} $0.70 = 0.577 + 0.00209\,(T - 600)$، إذن $T = 600 +
0.123/0.00209 \approx \qty{659}{\celsius}$.

\textbf{10.} \qty{7.0}{kbar}؛ و$\num{7.0e8}/101325 \approx \num{6900}$
ضغط جوي.

\textbf{11.} $\rho g = 2830 \times 9.81 \approx \qty{27800}{Pa}$ لكل متر،
أي \qty{27.8}{MPa} لكل كيلومتر.

\textbf{12.} $z = \num{0.70e9}/\num{27762} \approx \qty{25200}{m}$، أي نحو
\qty{25.2}{km}.

\textbf{13.} $\num{0.577e9}/\num{27762} \approx \qty{20.8}{km}$.

\textbf{14.} $\num{0.70e9}/(2700 \times 9.81) \approx \qty{26.4}{km}$، أي أعمق
بنحو \qty{5}{\percent} ($26.4/25.2 \approx 1.048$).

\textbf{15.} في القشرة: فالعمق \qty{25}{km} أقل حتى من سمك القشرة
العادي.

\textbf{16.} $(600 - 10)/25.2 \approx \qty{23.4}{\celsius}$ لكل كيلومتر.

\textbf{17.} إنه أدفأ بكثير من \omterm{def:g11:subduction-zones:subduction}{نطاق الاندساس} (دون
\qty{10}{\celsius} لكل كيلومتر) وأبرد قليلًا من القشرة النموذجية
(\qty{30}{\celsius})، وبعيد عن المنطقة البركانية (من 40 إلى \qty{50}{\celsius}): إنه
صخر دُفن عميقًا في القشرة المُسمَّكة لحزام جبلي، \omterm{def:g12:metamorphic-rocks:contact-regional}{بالتحول
الإقليمي}.

\textbf{18.} $10 + 30 \times 25.2 \approx \qty{766}{\celsius}$. وخط
الكيانيت--السيليمانيت، إذا مُدّ، يقع عند $0.577 + 0.00209 \times 166
\approx \qty{0.92}{GPa}$ عند تلك الدرجة؛ و\qty{0.70}{GPa} تحته،
في حقل السيليمانيت. \omterm{def:g11:lithosphere-asthenosphere:geotherm}{فالمنحنى الحراري الأرضي} النموذجي كان سيصنع السيليمانيت.

\textbf{19.} $\qty{25.2}{km} = \qty{2.52e7}{mm}$؛ وبمعدل \qty{0.5}{mm} في السنة،
نحو 50 مليون سنة.

\textbf{20.} كلما صعد الصخر برد، فتباطأت تفاعلاته،
والماء الذي تحتاج إليه التفاعلات العكسية كان قد طُرد في طريق
النزول.

\textbf{21.} تبلور شيست الكيانيت على عمق نحو \textbf{\qty{25}{km}
تحت السطح}، عند نحو \qty{600}{\celsius}.
\end{solution}
